\begin{abstract}
Embodied artificial intelligence (AI) systems are typically evaluated through benchmark pipelines that combine policies, checkpoints, simulators, task suites, reset states, interface adapters, seeds, and metric protocols. However, when an aggregate score changes, the available evidence is often insufficient to determine whether the deviation reflects a likely true regression, a setup-sensitive evaluation change, or an inconclusive outcome. This ambiguity limits the robustness and interpretability of embodied AI evaluation. This paper presents \sys, a factor-directed diagnosis framework for embodied AI evaluation deviations. \sys represents evaluation runs with structured manifests and episode-level outcomes, detects aggregate and paired-outcome deviations, and uses targeted differential replay to test factor-level hypotheses. It produces evidence-aware reports that classify deviation status, rank likely factors, and abstain when attribution is not sufficiently supported. We ground the problem through an empirical analysis of 473 GitHub issues and pull requests from 15 embodied AI and robot-learning projects, and evaluate \sys on a LeRobot--LIBERO artifact corpus covering completed rollouts, failed executions, calibration cases, ablations, robustness scenarios, and true-regression probes. \sys attributes all detected paper-only and robustness cases in our corpus to the expected factor, makes no false setup attributions on 12 negative calibration cases, and correctly classifies 66/67 main-status cases. Its factor-directed replay reduces diagnosis workload by 69.3\% compared with a blind rerun-$k$ baseline. These results show that embodied AI benchmarks can be made more reliable by complementing aggregate metrics with structured, factor-level diagnostic evidence.
\end{abstract}
